\documentclass[11pt,a4paper]{article}

\usepackage[utf8]{inputenc}
\usepackage[T1]{fontenc}
\usepackage[brazil]{babel}

\usepackage{amsmath}
\usepackage{amssymb}
\usepackage{graphicx}
\graphicspath{{figures/}}
\usepackage{booktabs}
\usepackage{float}
\usepackage{pgfplots}
\pgfplotsset{compat=1.18}
\usepackage{hyperref}
\usepackage{cleveref}

% Ajuste conforme o destino editorial.
% Para referências no padrão ABNT com biber:
% \usepackage[backend=biber,style=abnt]{biblatex}
% \addbibresource{references.bib}
% Caso use BibTeX clássico, remova as linhas acima e use \bibliography no final.

\title{Título do artigo}
\author{Autor A \and Autor B}
\date{}

\begin{document}

\maketitle

\begin{abstract}
Resumo do artigo: contexto, problema, objetivo, metodologia, principais
resultados e conclusão.
\end{abstract}

\noindent\textbf{Palavras-chave:} palavra1; palavra2; palavra3.

\input{sections/introduction}
\input{sections/literature}
\input{sections/methodology}
\input{sections/results}
\input{sections/discussion}
\input{sections/conclusion}

% Com biblatex/biber:
% \printbibliography

% Com BibTeX clássico:
% \bibliographystyle{plain}
% \bibliography{references}

\end{document}
